\begin{figure}

\begin{subfigure}{0.49\linewidth}
\centering
\includegraphics[width=\linewidth]{../binder/binder-2026-08-18-complexty-fitness-correlation.ipynb/binder/teeplots/2026-08-18-complexty-fitness-correlation/agg=True+viz=lmplot+when=all+x=genetic-complexity+y=focal-prevalence+ext=.pdf}

\vspace{-2ex}

\caption{\footnotesize genotype complexity, mean per replicate}
\label{fig:complexity-fitness-correlation-genetic-9s:genetic-agg}
\end{subfigure}
\begin{subfigure}{0.49\linewidth}
\centering
\includegraphics[width=\linewidth]{../binder/binder-2026-08-18-complexty-fitness-correlation.ipynb/binder/teeplots/2026-08-18-complexty-fitness-correlation/agg=False+viz=lmplot+when=all+x=genetic-complexity+y=focal-prevalence+ext=.pdf}

\vspace{-2ex}

\caption{\footnotesize genotype complexity, replicate samples taken at 10-stint intervals}
\label{fig:complexity-fitness-correlation-genetic-9s:genetic-unagg}
\end{subfigure}

\begin{subfigure}{0.49\linewidth}
\centering
\includegraphics[width=\linewidth]{../binder/binder-2026-08-18-complexty-fitness-correlation.ipynb/binder/teeplots/2026-08-18-complexty-fitness-correlation/agg=True+viz=lmplot+when=all+x=phenotype-complexity+y=focal-prevalence+ext=.pdf}

\vspace{-2ex}

\caption{\footnotesize phenotype complexity, mean per replicate}
\label{fig:complexity-fitness-correlation-genetic-9s:phenotype-agg}
\end{subfigure}
\begin{subfigure}{0.49\linewidth}
\centering
\includegraphics[width=\linewidth]{../binder/binder-2026-08-18-complexty-fitness-correlation.ipynb/binder/teeplots/2026-08-18-complexty-fitness-correlation/agg=False+viz=lmplot+when=all+x=phenotype-complexity+y=focal-prevalence+ext=.pdf}

\vspace{-2ex}

\caption{\footnotesize phenotype complexity, replicate samples taken at 10-stint intervals}
\label{fig:complexity-fitness-correlation-genetic-9s:phenotype-unagg}
\end{subfigure}

\vspace{-1ex}

% https://github.com/mmore500/oee4/blob/757995fa9b0b474395acd58e609ed6510ac9f4ca/binder/2026-08-18-complexty-fitness-correlation.ipynb
\caption{%
\textbf{Negligible correlation between fitness and complexity measures across 100-stint replicates, assayed at stints 9, 19, 29, 39, 49, 59, 69, 79, 89, 99.}
\footnotesize
Pearson correlation statistics annotated.
Shaded bands indicate bootstrap 95\% confidence for OLS regression.
Fitness values represent proportion ``measuring stick'' competitions won against comparator strains under no-background conditions.
Samples drawn from original 40-replicate/100-stint trials, which includes case study replicate.
Note that disaggregation artificially inflates statistical power in panels \labelcref{fig:complexity-fitness-correlation-genetic-9s:genetic-unagg,fig:complexity-fitness-correlation-genetic-9s:phenotype-unagg} due to introducing pseudoreplication; regardless, no relationship is detected in either case.
To maximize sensitivity, no corrections are made for multiple comparison.
Supplementary Figure \ref{fig:complexity-fitness-correlation-genetic} reports equivalent assays taken at stints 10, 20, 30, 40, 50, 60, 70, 80, 90, 100 --- which suffered sporadic assay job failures due to intermittent compute node/network failures arising with our service provider, but are included for completeness; regardless, a similar negligible relationship between fitness and genotype/phenotype complexity is detected in these trials.
(Due to practical considerations, it was operationally simpler and more relaible to cleanly collect new data at alternate timepoints after intermittent failure conditions had been remediated; this strategy also provided the opportunity to replicate original findings with independent --- i.e., offset by one stint --- specimen samples.)
Supplementary Figure \ref{fig:general-supp-nobg} analyzes fitness-complexity relationship in secondary 400-replicate/20-stint trials.
}
\label{fig:complexity-fitness-correlation-genetic-9s}

\end{figure}
